Este trabajo nace de un problema que el equipo enfrentaba de manera cotidiana en su propio entorno de trabajo: la creación de materiales en SAP consumía tiempo valioso en tareas de verificación que, en principio, podían automatizarse. Lo que comenzó como una limitante operativa se convirtió en la oportunidad de aplicar, sobre un problema real y medible, las herramientas de inteligencia de negocios y ciencia de datos desarrolladas a lo largo de la maestría.

Uno de los retos particulares de este trabajo fue mantener el equilibrio entre la confidencialidad de la información utilizada y el rigor académico exigido. La organización donde se desarrolló e implementó el proyecto autorizó el uso de su información bajo la condición de que no se revelara su identidad; por esa razón, a lo largo de este documento se hace referencia a ella de forma genérica, y se omiten los nombres de las personas involucradas en el proceso operativo.

Otra limitación reconocida es el alcance acotado del proyecto: los resultados presentados corresponden a una única unidad de negocio y a un período de validación específico, por lo que deben interpretarse dentro de ese contexto y no como una generalización automática a otros entornos.

Finalmente, este documento se benefició de que el proyecto no fue un ejercicio teórico, sino una herramienta que estuvo en uso activo por personas reales durante varias semanas. Esa validación con datos de producción y con retroalimentación de los usuarios finales es lo que le da a este trabajo su mayor valor.